% P16.17: consistency diagnosis, corrected hitting-set search and measurements.
\chapter{Consistency-Based Diagnosis}
\label{sec:p16-consistency-based-diagnosis}

A system model specifies how components behave when they function
normally. An observation can contradict the conjunction of those
normality assumptions. Diagnosis asks which assumptions must be
withdrawn to make the model and observation jointly possible.
The input fixes the component list and the system description.
An omitted component, a faulty observation or an incorrect system
description therefore changes the diagnostic problem.

We develop Reiter's model, its conflict characterization, measurement
theory and default-logic correspondence
\cite{Reiter1987Diagnosis}. The search construction includes the
shared-node and relabeling rules of Greiner, Smith and Wilkerson
\cite[secs.~3--4, author version]{GreinerSmithWilkerson1989Correction}.
All uses of minimality below mean inclusion minimality unless a
cardinality or probability criterion is explicitly stated.

\section{Normality, exact fault assignments and withdrawal}
\label{sec:p16-diagnosis-semantics}

\begin{hbdefinition}{def:p16-diagnosis-input}{A diagnosis problem}
Fix a language with a class of structures and bivalent satisfaction
for its sentences. Negation, conjunction and disjunction have their
classical truth conditions. A theory is \emph{consistent} when it
has a structure in that class; semantic consequence is written
\(\entails\).

Let \(C\) be a finite set of component names and let
\(a_c\), for \(c\in C\), be designated abnormality sentences.
In Reiter's first-order presentation \(a_c\) is
\(\mathrm{AB}(c)\). Let \(S\) be the system description and \(O\)
the observations, both sets of sentences, and put \(W=S\cup O\).
For \(N,\Delta\subseteq C\), define
\[
\begin{split}
 \mathsf{Normal}(N)&=\{\neg a_c:c\in N\},\\
 W^-_\Delta&=W\cup\mathsf{Normal}(C\setminus\Delta),\\
 W^=_\Delta&=W^-_\Delta\cup\{a_c:c\in\Delta\}.
\end{split}
\]
A \emph{withdrawal set} is a \(\Delta\) for which \(W^-_\Delta\)
is consistent. An \emph{exact fault assignment} is a \(\Delta\)
for which \(W^=_\Delta\) is consistent. A \emph{diagnosis} is
an inclusion-minimal exact fault assignment. Write
\(\mathsf D(W)\) for the diagnosis family.
No independence, probability or prescribed behavior of faulty
components is part of this definition. The theory may relate
abnormalities, for example by \(a_c\to a_d\).
\index{diagnosis!consistency based}\index{diagnosis!normality assumption}
\end{hbdefinition}

\begin{hbtheorem}{thm:p16-diag-semantics}{proved=annex:p16.diag-semantics}{The withdrawal characterization and all empty cases}
For the preceding input:
\begin{enumerate}
\item The diagnoses are precisely the inclusion-minimal withdrawal
      sets. For every diagnosis \(\Delta\),
      \(W^-_\Delta\entails a_c\) for each \(c\in\Delta\);
      consequently \(W^-_\Delta\) and \(W^=_\Delta\) have the
      same structures.
\item A diagnosis exists exactly when \(W\) is consistent.
      The empty set is a diagnosis exactly when
      \(W\cup\mathsf{Normal}(C)\) is consistent; in that case
      it is the sole diagnosis.
\item Withdrawal sets are closed under enlargement, and every
      withdrawal set contains a diagnosis. Exact fault assignments
      need not be closed under enlargement.
\item To test that a withdrawal set \(\Delta\) is minimal, it is
      sufficient and necessary to test that
      \(W^-_{\Delta\setminus\{c\}}\) is inconsistent for every
      \(c\in\Delta\).
\end{enumerate}
These statements include \(C=\varnothing\) and inconsistent \(W\).
\end{hbtheorem}

For example, with \(C=\{c,d\}\) and \(W=\{a_c\to a_d\}\),
\(\varnothing\) is the sole diagnosis. Withdrawing \(c\)'s normality
is consistent, but declaring exactly \(c\) abnormal is inconsistent.
The equivalence in the theorem concerns the minimal sets.

\section{Conflicts and hitting sets}
\label{sec:p16-diagnosis-conflicts}

\begin{hbdefinition}{def:p16-diagnosis-conflicts}{Conflicts and transversals}
A \emph{conflict} is a \(K\subseteq C\) such that
\(W\cup\mathsf{Normal}(K)\) is inconsistent. Let
\(\mathcal K(W)\) be all conflicts and
\(\mathcal K_{\min}(W)\) its inclusion-minimal members.
For any family \(\mathcal F\subseteq\mathcal P(C)\), a
\emph{hitting set} is an \(H\subseteq C\) satisfying
\(H\cap K\ne\varnothing\) for every \(K\in\mathcal F\).
Let \(\mathsf{Tr}(\mathcal F)\) be its minimal hitting sets.
The empty family has the single minimal hitting set \(\varnothing\).
A family containing \(\varnothing\) has no hitting set.
\index{diagnosis!conflict}\index{hitting set!minimal}
\end{hbdefinition}

\begin{hbtheorem}{thm:p16-diag-conflicts}{proved=annex:p16.diag-conflicts}{Conflict duality}
Conflicts are closed under enlargement and
\[
\begin{split}
 W^-_\Delta\text{ is consistent}
   &\quad\Longleftrightarrow\quad
        \Delta\text{ hits every member of }\mathcal K(W),\\
 \mathsf D(W)
   &=\mathsf{Tr}(\mathcal K(W))
     =\mathsf{Tr}(\mathcal K_{\min}(W)).
\end{split}
\]
Minimality of a hitting set \(H\) is equivalent to the existence,
for every \(c\in H\), of a conflict \(K_c\) with
\(K_c\cap H=\{c\}\).
For every finite family \(\mathcal F\subseteq\mathcal P(C)\),
\[
 \mathsf{Tr}(\mathsf{Tr}(\mathcal F))
       =\min_{\subseteq}\mathcal F.
\]
Here \(\min_{\subseteq}\) denotes the family of inclusion-minimal
members. Both exceptional empty-family cases obey this identity.
\end{hbtheorem}

\begin{hbtheorem}{thm:p16-diag-singletons}{proved=annex:p16.diag-singletons}{Single faults and partial conflict information}
Suppose at least one conflict exists. Then
\[
 \{c\}\in\mathsf D(W)
   \quad\Longleftrightarrow\quad
 c\in\bigcap\mathcal K(W)
   \quad\Longleftrightarrow\quad
 c\in\bigcap\mathcal K_{\min}(W).
\]
If \(K_1,\ldots,K_r\), \(r\ge1\), are any known conflicts, then
\[
 \{c\}\in\mathsf D(W)
 \quad\Longleftrightarrow\quad
 c\in\bigcap_{i=1}^rK_i
 \quad\hbox{and}\quad W^-_{\{c\}}\text{ is consistent}.
\]
Without a conflict, the empty diagnosis is already the sole
diagnosis. Intersecting an empty collection of conflicts must not
be used to declare every component a single fault.
\end{hbtheorem}

\section{Complete search with explicit consistency access}
\label{sec:p16-diagnosis-search}

\begin{hbdefinition}{def:p16-diagnosis-hs-tree}{An unpruned hitting-set tree}
Let \(C\) be finite and \(\mathcal F\subseteq\mathcal P(C)\).
Each node carries the set \(H\) of edge labels on its root path.
The root has \(H=\varnothing\). At a node, choose a
\(K\in\mathcal F\) disjoint from \(H\), if one exists.
Label the node \(K\), and create an edge labeled \(c\) to a child
with path set \(H\cup\{c\}\), for each \(c\in K\).
If no such \(K\) exists, label the node \(\checkmark\) and give
it no children. A node labeled by the empty conflict has no
children and is not a success node.
\end{hbdefinition}

\begin{hbtheorem}{thm:p16-diag-hs-tree}{proved=annex:p16.diag-hs-tree}{Completeness of the basic tree}
Every branch has length at most \(|C|\). Every success path
is a hitting set, and every minimal hitting set is a success
path set. Success path sets can be nonminimal; taking their
inclusion-minimal members returns exactly
\(\mathsf{Tr}(\mathcal F)\).
The conclusion does not require minimal conflicts or a special
choice order for the disjoint conflicts.
\end{hbtheorem}

\begin{hbdefinition}{def:p16-diagnosis-hs-dag}{A corrected shared-node search}
The following finite procedure works with an explicitly listed
family \(\mathcal F\). The oracle construction below supplies
its implicit diagnostic version.
\begin{enumerate}
\item Store a directed acyclic graph with a unique active node
      for each represented set \(H\subseteq C\). The root is
      \(\varnothing\). Process active unprocessed nodes in
      increasing order of \(|H|\), breaking ties by a fixed order.
\item If a success node \(J\) satisfies \(J\subsetneq H\),
      close \(H\), without a label query or successors.
      Otherwise find a member \(K\) of the current family
      disjoint from \(H\). If there is none, mark \(H\)
      successful. Otherwise label it \(K\).
\item For every \(c\in K\), connect \(H\) to the node
      \(H\cup\{c\}\). Reuse that node if it already exists.
      Preserve all of its incoming edges; do not replace a
      second incoming edge by an unconnected closed copy.
\item When a newly used conflict \(K\) is a proper subset of
      an earlier label \(K'\), relabel every such node by \(K\).
      At that node remove precisely the outgoing edges whose
      labels belong to \(K'\setminus K\). Delete a node only
      if it is no longer reachable from the root. A surviving
      node with another incoming path keeps its descendants.
      If this operation removes the node being processed,
      discontinue processing that node.
\item Remove each superseded \(K'\) from the stored family.
      Repeat the preceding relabeling for all superseded
      labels, then continue processing reachable nodes.
      Return the success sets when no unprocessed node remains.
\end{enumerate}
Removing a superseded set implements its redundancy in the
family. The source also describes placing its subset ahead of
it in the ordered conflict list; that gives the same available
disjoint-conflict test. Graph shapes and the order of equally
large answers can differ.
\index{diagnosis!hitting-set DAG}
\end{hbdefinition}

\begin{hbtheorem}{thm:p16-diag-hs-dag}{proved=annex:p16.diag-hs-dag}{Correctness of shared-node closing and pruning}
The procedure terminates and returns exactly
\(\mathsf{Tr}(\mathcal F)\), in nondecreasing cardinality.
Every active edge adds one element to its source set, so
all paths to a shared node have the same set and length.
The graph has at most \(2^{|C|}\) active nodes and
\(|C|2^{|C|}\) active edges.
The relabeling operation preserves the disjoint-conflict
condition and every successor required by the smaller label.
Both node sharing and retention of nodes reachable through
another parent are required for the stated pruning procedure.
\end{hbtheorem}

\begin{hbdefinition}{def:p16-diagnosis-oracle}{The diagnostic consistency oracle}
A \emph{total diagnostic oracle} takes \(N\subseteq C\) and
returns either a conflict \(K\subseteq N\), or
\(\checkmark\) certifying that \(W\cup\mathsf{Normal}(N)\)
is consistent. It must always terminate, and every answer must
have its stated meaning. A conflict need not be minimal.
\end{hbdefinition}

\begin{hbtheorem}{thm:p16-diag-oracle}{proved=annex:p16.diag-oracle}{Implicit conflict search and its computational premise}
In the shared-node procedure, keep a cache of previously
returned conflicts. Before calling the oracle at \(H\), reuse
any cached conflict disjoint from \(H\). If none exists, call
the oracle with \(C\setminus H\). Apply the same relabeling,
node-sharing and closing operations to each new conflict.
When a cached conflict is removed as redundant, retain its
smaller replacing conflict in the cache.

This procedure terminates and returns all and only the
diagnoses. No two oracle calls return the same conflict;
a later returned conflict never contains an earlier returned
conflict. Stopping after all nodes of cardinality at most
\(k\) have been processed supplies all diagnoses of
cardinality at most \(k\).

For a finite propositional theory \(W\) whose sentences and
designated \(a_c\)'s use a total of \(v\) atoms, a total
oracle exists: enumerate the \(2^v\) valuations. If no
valuation satisfies the current normality assumptions,
return the whole queried set \(N\); otherwise return
\(\checkmark\). At most \(|N|\) further consistency tests
suffice to shrink a returned conflict to a minimal one.
For an arbitrary first-order theory the total-oracle premise
is additional; a refutation procedure alone need not terminate
on a consistent query.
\end{hbtheorem}

\begin{hbtheorem}{thm:p16-diag-pruning-calculations}{proved=annex:p16.diag-pruning-calculations}{Worked hitting sets and retained paths}
\begin{enumerate}
\item For
 \[
 \mathcal F=\{\{2,4,5\},\{1,2,3\},\{1,3,5\},
                 \{2,4,6\},\{2,4\},\{2,3,5\},\{1,6\}\},
 \]
 the minimal hitting sets are exactly
 \[
 \{1,2\},\quad\{2,3,6\},\quad\{2,5,6\},\quad
 \{1,3,4\},\quad\{1,4,5\},\quad\{3,4,6\}.
 \]
\item For
 \(\mathcal G=\{\{a,b\},\{b,c\},\{a,c\},\{b,d\},\{b\}\}\),
 the answers are \(\{a,b\}\) and \(\{b,c\}\).
 After the root is relabeled \(\{b\}\), the shared node
 \(\{a,b\}\) remains reachable through the root's \(b\)-edge
 and then the \(a\)-edge. Deleting its other incoming branch
 therefore preserves that answer.
\item For \(\{\{a,b\},\{a\},\{b\}\}\), the sole answer is
 \(\{a,b\}\). Relabeling by \(\{b\}\) removes the root's
 \(a\)-edge but does not subsequently license deletion of
 its \(b\)-edge on comparison with the incomparable
 conflict \(\{a\}\).
\end{enumerate}
\end{hbtheorem}

\section{Two fully specified circuit calculations}
\label{sec:p16-diagnosis-circuits}

\begin{hbdefinition}{def:p16-diagnosis-circuit-inputs}{Circuit variables and guarded equations}
For a Boolean full adder, let
\[
 C=\{X_1,X_2,A_1,A_2,O_1\},\qquad x,u,v,y,z\in\{0,1\}.
\]
The inputs are \(1,0,1\); the observed outputs are \(y=1,z=0\).
The system requires the following equations when the named
component is normal:
\[
\begin{array}{c|ccccc}
 \text{component}&X_1&A_1&A_2&X_2&O_1\\ \hline
 \text{normal equation}&x=1&u=0&v=x&y=1-x&z=u\lor v .
\end{array}
\]
The symbol \(\lor\) in the last equation is Boolean disjunction.
All wire values remain Boolean even when a component is faulty.

For the multiplier-adder network let
\[
 C=\{M_1,M_2,M_3,A_1,A_2\},\qquad x_1,x_2,x_3\in\mathbb Z .
\]
Each multiplier has inputs \(3,2\), and the observed adder
outputs are \(10,12\). Normal \(M_i\) requires \(x_i=6\);
normal \(A_1\) requires \(x_1+x_2=10\); normal \(A_2\)
requires \(x_2+x_3=12\). An abnormal component contributes
no additional equation. These are stipulated mathematical
systems, with observations already incorporated in \(W\).
\end{hbdefinition}

\begin{hbtheorem}{thm:p16-diag-circuits}{proved=annex:p16.diag-circuits}{All diagnoses of the two circuits}
The Boolean adder has precisely two minimal conflicts,
\[
             \{X_1,X_2\},\qquad \{X_1,A_2,O_1\},
\]
and precisely three diagnoses,
\[
                 \{X_1\},\qquad\{X_2,A_2\},\qquad\{X_2,O_1\}.
\]
It has ten conflicts in total.
The multiplier-adder network has precisely two minimal conflicts,
\[
           \{M_1,M_2,A_1\},\qquad\{M_1,M_3,A_1,A_2\},
\]
and precisely four diagnoses,
\[
                 \{M_1\},\quad\{A_1\},\quad
                 \{M_2,M_3\},\quad\{M_2,A_2\}.
\]
Each listed diagnosis has an explicit consistent wire assignment.
\end{hbtheorem}

\section{What new measurements change}
\label{sec:p16-diagnosis-measurements}

\begin{hbdefinition}{def:p16-diagnosis-prediction}{Prediction and a measurement update}
A diagnosis \(\Delta\) \emph{predicts} a sentence \(\pi\)
when \(W^=_\Delta\entails\pi\), equivalently
\(W^-_\Delta\entails\pi\). A measurement update adds a
sentence \(\pi\), producing \(W'=W\cup\{\pi\}\).
A finite conjunction represents several new measurements.
Failure to predict \(\pi\) does not imply prediction of
\(\neg\pi\).
\index{diagnosis!measurement}\index{diagnosis!prediction}
\end{hbdefinition}

\begin{hbtheorem}{thm:p16-diag-measurements}{proved=annex:p16.diag-measurements}{Persistence, rejection and new minimal diagnoses}
\begin{enumerate}
\item An old diagnosis \(\Delta\) remains a diagnosis after
      measuring \(\pi\) exactly when it does not predict
      \(\neg\pi\). In particular confirming measurements
      preserve it, and disconfirming measurements reject it.
\item Every new diagnosis contains an old diagnosis.
      If it is not itself old, every old diagnosis that
      it contains is a proper subset and predicts
      \(\neg\pi\).
\item If no old diagnosis predicts \(\neg\pi\), the two
      diagnosis families are identical.
\item If the empty set was not an old diagnosis, every
      diagnosis that is new has at least two components.
      Hence new single-fault diagnoses cannot arise.
      If every old diagnosis decides \(\pi\), the surviving
      single faults are precisely those predicting \(\pi\).
\end{enumerate}
The first three statements also cover an inconsistent initial
or updated theory. A measurement can change predictions even
when it leaves the diagnosis family fixed.
\end{hbtheorem}

\begin{hbtheorem}{thm:p16-diag-measurement-calculations}{proved=annex:p16.diag-measurement-calculations}{Complete measurement updates for the multiplier network}
In the multiplier-adder network:
\begin{enumerate}
\item Diagnoses \(\{M_1\}\) and \(\{A_1\}\) predict \(x_2=6\).
      Diagnoses \(\{M_2,M_3\}\) and \(\{M_2,A_2\}\)
      predict \(x_2=4\).
\item After measuring \(x_2=5\), the diagnoses are precisely
 \[
 \{M_1,M_2,M_3\},\quad\{M_1,M_2,A_2\},\quad
 \{A_1,M_2,M_3\},\quad\{A_1,M_2,A_2\}.
 \]
      Every old diagnosis is rejected.
\item After measuring \(x_1=6\) instead, the diagnoses are
      precisely \(\{A_1\},\{M_2,M_3\},\{M_2,A_2\}\).
      Thus \(\{A_1\}\) is the only remaining single fault,
      although multiple-fault diagnoses remain.
\end{enumerate}
\end{hbtheorem}

\section{Defaults, covering models and explanatory entailment}
\label{sec:p16-diagnosis-defaults}

\begin{hbdefinition}{def:p16-diagnosis-defaults}{Normality defaults}
In classical first-order logic, let \(\operatorname{Th}\)
denote deductive closure and form the closed default theory
\[
       \mathcal T_W=
       \left(\left\{\frac{\top:\neg a_c}{\neg a_c}:c\in C\right\},W\right).
\]
Use Reiter extensions as defined in
\cref{def:p18-default-closure}. For this prerequisite-free
theory the operator on a closed candidate \(E\) is
\[
 \Gamma(E)=\operatorname{Th}
       \bigl(W\cup\{\neg a_c:c\in C,\ a_c\notin E\}\bigr).
\]
An extension satisfies \(E=\Gamma(E)\).
\end{hbdefinition}

\begin{hbtheorem}{thm:p16-diag-defaults}{proved=annex:p16.diag-defaults}{The complete default correspondence and its consistency condition}
If \(W\) is consistent, the map
\[
       \Delta\longmapsto\operatorname{Th}(W^-_\Delta)
\]
is a bijection from diagnoses to Reiter extensions of
\(\mathcal T_W\). Its inverse is
\(E\mapsto\{c\in C:a_c\in E\}\). Each extension consists
exactly of the sentences predicted by its diagnosis.

More generally, for finitely many closed sentences
\(\alpha_i\) and consistent \(W\), the extensions of
the default theory \((\{\top:\alpha_i/\alpha_i\},W)\) are exactly
\(\operatorname{Th}(W\cup A)\), where \(A\) is a maximal
consistent-with-\(W\) subset of \(\{\alpha_i\}\).
If \(W\) is inconsistent, the normality default theory
has the sole inconsistent extension, the whole sentence
language, while the diagnosis family is empty.
\end{hbtheorem}

\begin{hbdefinition}{def:p16-diagnosis-cover-model}{A covering input and its logical representation}
Let \(D\) be a finite disorder set, \(M\) a finite manifestation
set, \(R\subseteq D\times M\), and \(M^+\subseteq M\).
Put \(R_m=\{d\in D:(d,m)\in R\}\). A cover is
\(\Delta\subseteq D\) such that \(\Delta\cap R_m\ne\varnothing\)
for every \(m\in M^+\).

Introduce separate indexed atoms \(a_d,p_d,o_m\).
Take \(C=D\) and the theory
\[
 S=\{\neg a_d\to\neg p_d:d\in D\}
       \cup\{o_m\to\bigvee_{d\in R_m}p_d:m\in M\},
       \qquad O=\{o_m:m\in M^+\}.
\]
An empty disjunction is false. Indexing keeps the three
atom families disjoint even if \(D\) and \(M\) overlap
as underlying sets. The interpretation of \(R\) as
possible causal production is supplied by an application;
the displayed theory itself is a Boolean constraint model.

Its named first-order version uses distinct indexed names for
disorders and manifestations and unary predicates
\(\mathrm{Dis},\mathrm{Present},\mathrm{Observed},\mathrm{AB}\).
The system axioms are \(\mathrm{Dis}(d)\) for \(d\in D\),
\[
 \mathrm{Observed}(m)\to
       \bigvee_{d\in R_m}\mathrm{Present}(d)\quad(m\in M),
 \qquad
 \forall x\,(\mathrm{Dis}(x)\land\neg\mathrm{AB}(x)
                         \to\neg\mathrm{Present}(x)).
\]
The observations are \(\mathrm{Observed}(m)\) for
\(m\in M^+\), and \(a_d=\mathrm{AB}(d)\).
The annex gives structures in both directions, including
empty disorder or manifestation sets.
\index{diagnosis!set covering}\index{diagnosis!default logic}
\end{hbdefinition}

\begin{hbtheorem}{thm:p16-diag-set-cover}{proved=annex:p16.diag-set-cover}{Covering and diagnosis}
For this representation, withdrawal sets are precisely
the covers of \(M^+\), and diagnoses are precisely its
inclusion-minimal covers. The minimal conflicts are
the inclusion-minimal members of
\(\{R_m:m\in M^+\}\).
If an observed manifestation has \(R_m=\varnothing\),
there is no diagnosis. If \(M^+=\varnothing\), the sole
diagnosis is empty.

Minimum-cardinality covers form a subfamily of the diagnoses,
and the inclusion can be strict. For
\[
 R_{m_1}=\{a,b\},\qquad R_{m_2}=\{a,c\},\qquad
                 M^+=\{m_1,m_2\},
\]
the diagnoses are \(\{a\}\) and \(\{b,c\}\), while only
\(\{a\}\) has minimum cardinality.
\end{hbtheorem}

\begin{hbtheorem}{thm:p16-diag-abduction-boundaries}{proved=annex:p16.diag-abduction-boundaries}{Consistency, entailment and declared fault modes}
\begin{enumerate}
\item With one component, atom \(o\), system
      \(S=\{\neg a_c\to\neg o\}\), and observation \(o\),
      the sole diagnosis is \(\{c\}\). Nevertheless
      \(S\cup\{a_c\}\not\entails o\).
\item Replacing \(S\) by \(S\cup\{a_c\to o\}\) preserves
      that diagnosis and makes \(S\cup\{a_c\}\entails o\)
      for the augmented \(S\).
      The extra fault-mode implication supplies the
      explanatory entailment.
\item With \(C=\{c,d\}\) and
      \(W=\{a_c\lor a_d,\neg(a_c\land a_d)\}\),
      the diagnoses are \(\{c\}\) and \(\{d\}\).
      One has \(W\not\entails a_c\) and \(W\not\entails a_d\).
      The superset \(\{c,d\}\) is a withdrawal
      set but not an exact fault assignment.
\end{enumerate}
Thus the consistency requirement of diagnosis and the
observation-entailment requirement of
\cref{def:p16-abduction-models} impose different conditions.
\end{hbtheorem}

\begin{hbtheorem}{thm:p16-diag-probability-boundary}{proved=annex:p16.diag-probability-boundary}{What an independent fault prior ranks}
Let independent fault indicators have probabilities
\(0<p_c<1\). Write \(P(\Delta)\) for the probability of
the event that exactly the indicators indexed by
\(\Delta\subseteq C\) equal one. Then
\[
 P(\Delta)=
   \prod_{c\in\Delta}p_c\prod_{c\notin\Delta}(1-p_c)
   =\left(\prod_{c\in C}(1-p_c)\right)
        \prod_{c\in\Delta}\frac{p_c}{1-p_c}.
\]
Every added component lowers this probability if all
\(p_c<1/2\), but among incomparable fault sets fewer
components need not mean higher probability.
For \(p_a=1/10000\) and \(p_b=p_c=1/10\),
the two-fault set \(\{b,c\}\) is more probable than
the singleton \(\{a\}\).
If all probabilities equal one common \(p<1/2\),
probability decreases strictly with cardinality.

For an observed event \(E\) with positive probability, posterior
comparison also multiplies each prior by its likelihood
\(P(E\mid\Delta)\). Logical consistency alone does not
specify those likelihoods.
\end{hbtheorem}

The results concern the stated finite component universe,
including simultaneous faults and interactions represented in
the theory. They do not locate failures outside that universe
or justify a probability model for a physical system.
The separate original accounts of probabilistic diagnosis,
fault-mode abduction and protected-premise counterfactual
diagnosis require their own input semantics and derivations.
